% !TEX program = xelatex
\documentclass[11pt]{article}
\usepackage[margin=0.75in]{geometry}
\usepackage{parskip}
\usepackage{fontspec}
\usepackage{xcolor}
\usepackage{titlesec}
\usepackage{enumitem}
\usepackage{tcolorbox}
\usepackage{multicol}
\usepackage{graphicx}
\usepackage{hanging}

% Fonte padrão do Windows (Arial)
\setmainfont{Arial}[
  Ligatures=TeX,
  Scale=MatchLowercase
]

% Cores SPE
\definecolor{spered}{RGB}{230, 0, 0}
\definecolor{spedark}{RGB}{0, 40, 80}
\definecolor{lightgray}{RGB}{245, 245, 245}
\definecolor{boxblue}{RGB}{220, 240, 255}
\definecolor{boxgreen}{RGB}{230, 255, 230}
\definecolor{boxyellow}{RGB}{255, 255, 200}
\definecolor{boxorange}{RGB}{255, 235, 200}
\definecolor{boxpurple}{RGB}{240, 220, 255}

% Configuração de títulos
\titleformat{\section}{\large\bfseries\color{spedark}}{\thesection}{1em}{}
\titlespacing*{\section}{0pt}{10pt}{5pt}

% Configuração base do tcolorbox
\tcbset{
  colback=white,
  colframe=spedark,
  arc=2mm,
  boxrule=0.5pt,
  left=6pt,
  right=6pt,
  top=6pt,
  bottom=6pt,
  fonttitle=\bfseries\small\centering,  % Título centrado
  coltitle=white,
  colbacktitle=spedark,
  titlerule=0pt,                         % Remove linha extra se houver
  before title={\hspace{0pt}},            % Ajuste fino
}

% Listas compactas
\setlist{nosep, leftmargin=*, before=\vspace{2pt}, after=\vspace{2pt}, itemsep=2pt}

% Cabeçalho
\title{\vspace{-1cm}\Huge\bfseries\color{spedark} Plano de Atividades \\[0.2cm] \Large Capítulo Estudantil SPE do ISPTEC}
\author{}
\date{}

\begin{document}
\maketitle

\vspace{-0.8cm}

\begin{center}
  \small\textcolor{spedark}{\textbf{Propostas para 2025 -- baseadas em boas práticas de capítulos SPE internacionais e adaptadas à realidade angolana.}}
\end{center}

\vspace{0.3cm}

\begin{multicols}{2}

% ==================== COLUNA 1 ====================

\tcbset{colback=boxblue}
\begin{tcolorbox}[title={\large\color{spedark}📚 Educacionais e Técnicos}]
  \begin{itemize}
    \item \textbf{Workshops temáticos mensais:} Exploração de petróleo, produção, refino, sustentabilidade, IA aplicada. Ministrados por estudantes finalistas ou docentes.
    \item \textbf{Ciclo de palestras “SPE Talks”:} Convidar engenheiros da Sonangol, Chevron, TotalEnergies para partilhar experiências e tendências do setor.
    \item \textbf{Visitas técnicas:} Programar idas à Refinaria de Luanda, campos onshore (ex.: Cabo Ledo) e estaleiros navais.
    \item \textbf{PetroBowl interno:} Competição de perguntas e respostas sobre engenharia petrolífera, estilo quiz, com prémios simbólicos.
    \item \textbf{Sessões de software técnico:} Introdução ao Petrel, Eclipse, Python para modelagem de reservatórios.
    \item \textbf{Clube de leitura de papers SPE:} Encontros quinzenais para discutir artigos científicos atuais.
  \end{itemize}
\end{tcolorbox}

\tcbset{colback=boxgreen}
\begin{tcolorbox}[title={\large\color{spedark}🤝 Networking e Integração}]
  \begin{itemize}
    \item \textbf{Boas-vindas semestral:} Icebreakers, apresentação do board, coffee break com música angolana.
    \item \textbf{Encontro com alumni:} Profissionais formados no ISPTEC partilham percursos e dicas.
    \item \textbf{Parceria com outros capítulos:} Eventos conjuntos com UAN, UEA, e internacionalmente com Brasil e Nigéria (webinars).
    \item \textbf{Mentoria entre pares:} Alunos mais velhos acompanham caloiros em grupos pequenos.
    \item \textbf{Networking dinner:} Jantar temático com profissionais da indústria em Luanda.
    \item \textbf{LinkedIn boost:} Sessão prática de otimização de perfis e conexões estratégicas.
  \end{itemize}
\end{tcolorbox}

\tcbset{colback=boxyellow}
\begin{tcolorbox}[title={\large\color{spedark}💰 Captação de Recursos}]
  \begin{itemize}
    \item \textbf{Mensalidades simbólicas:} 5.000–10.000 Kz/ano por membro ativo.
    \item \textbf{Patrocínios empresariais:} Propostas para Sonangol, TotalEnergies, Chevron, BAI.
    \item \textbf{Grants da SPE:} Candidaturas anuais a \emph{Student Chapter Grants} (até USD 1.000).
    \item \textbf{Venda de merchandise:} Camisetas, canecas, pins com logótipo do capítulo.
    \item \textbf{Crowdfunding:} Campanhas na plataforma Kuenda para projetos específicos (ex.: laboratório virtual).
    \item \textbf{Feijoada solidária:} Evento gastronómico aberto à comunidade académica.
  \end{itemize}
\end{tcolorbox}

% ==================== COLUNA 2 ====================
\columnbreak

\tcbset{colback=boxorange}
\begin{tcolorbox}[title={\large\color{spedark}🌱 Social e Ambiental}]
  \begin{itemize}
    \item \textbf{Limpeza de praias:} Ação de voluntariado na praia de Luanda, em parceria com ONGs locais.
    \item \textbf{Oficinas de energia solar:} Ensinar comunidades do Bengo a montar pequenos painéis solares.
    \item \textbf{Campanha de reciclagem:} Recolha de plástico no campus, venda para recicladores.
    \item \textbf{Dia da Sustentabilidade:} Palestras e exposições sobre ESG e transição energética em Angola.
    \item \textbf{Plantar árvores:} Parceria com o governo provincial para reflorestamento urbano.
    \item \textbf{Palestras em escolas:} Levar noções de energia e ambiente a estudantes do ensino médio.
  \end{itemize}
\end{tcolorbox}

\tcbset{colback=boxpurple}
\begin{tcolorbox}[title={\large\color{spedark}🚀 Desenvolvimento de Carreira}]
  \begin{itemize}
    \item \textbf{Feira de estágios:} Convidar empresas para apresentar vagas e recrutar no campus.
    \item \textbf{Workshop de CV e entrevistas:} Dicas práticas com profissionais de RH da indústria.
    \item \textbf{Programa de mentoria profissional:} Engenheiros experientes acompanham alunos finalistas.
    \item \textbf{Simulação de processos seletivos:} Dinâmicas de grupo e cases reais.
    \item \textbf{Certificações SPE:} Apoio na preparação para o exame de certificação (PEC).
    \item \textbf{Trilhas de carreira:} Painéis com profissionais de diferentes áreas (upstream, midstream, downstream, renováveis).
  \end{itemize}
\end{tcolorbox}

\tcbset{colback=lightgray}
\begin{tcolorbox}[title={\large\color{spedark}📱 Digitais e Inovação}]
  \begin{itemize}
    \item \textbf{Podcast “PetroCast”:} Entrevistas com profissionais e professores, disponível no Spotify.
    \item \textbf{Desafios online:} Quizzes sobre energia no Instagram e TikTok, com prémios.
    \item \textbf{Visitas virtuais:} Utilizar vídeos 360° ou YouTube para mostrar instalações petrolíferas.
    \item \textbf{Plataforma de cursos:} Compilar materiais gratuitos (SPE, Coursera, Petrobras) num Google Classroom.
    \item \textbf{Aplicativo do capítulo:} Desenvolvido por alunos de Informática com agenda, recursos e fórum.
    \item \textbf{Hackathon de energia:} Competição de ideias para soluções energéticas em Angola.
  \end{itemize}
\end{tcolorbox}

\end{multicols}

\vspace{0.3cm}

% Calendário semestral detalhado


\vspace{0.2cm}

\begin{center}
  \footnotesize\textcolor{gray}{\textit{Ideias recolhidas de capítulos SPE da Universidade de Houston, Instituto Superior Técnico (Portugal), Universidade de Lagos (Nigéria) e adaptadas ao contexto do ISPTEC.}}
\end{center}

\vspace{0.1cm}

\begin{center}
  \footnotesize\textcolor{spedark}{\textbf{Próximos passos:} O board irá priorizar as atividades com base nos recursos disponíveis e feedback dos membros.}
\end{center}

\begin{center}
  \footnotesize\textcolor{gray}{\textbf{Feito por:} Rocélio Da Silva.}
\end{center}

\end{document}